\documentclass[11pt,a4paper]{article}

\usepackage[T1]{fontenc}
\usepackage[utf8]{inputenc}
\usepackage{lmodern}
\usepackage{microtype}
\usepackage{amsmath,amssymb,amsthm}
\usepackage{geometry}
\usepackage{hyperref}
\usepackage{xcolor}
\usepackage{listings}
\usepackage{enumitem}
\usepackage{parskip}
\usepackage{titlesec}
\usepackage{booktabs}
\usepackage{fancyhdr}
\usepackage{abstract}

\geometry{
  a4paper,
  top=2cm, bottom=2cm,
  left=3cm, right=3cm
}
\emergencystretch=2em

\definecolor{codeblue}{RGB}{20,115,127}
\definecolor{codegray}{RGB}{245,245,245}
\definecolor{codered}{RGB}{160,0,0}
\definecolor{linkblue}{RGB}{10,60,150}

\hypersetup{
  colorlinks=true,
  linkcolor=linkblue,
  citecolor=linkblue,
  urlcolor=linkblue,
  pdftitle={Cascading Discreet Log Contracts},
  pdfauthor={Pedro Martins Rodrigues Novaes}
}

\lstset{
  backgroundcolor=\color{codegray},
  basicstyle=\ttfamily\small,
  breaklines=true,
  breakatwhitespace=false,
  frame=none,
  keepspaces=true,
  showstringspaces=false,
  tabsize=2,
  xleftmargin=1em,
  xrightmargin=1em,
  aboveskip=0.7em,
  belowskip=0.7em,
}

\newcommand{\code}[1]{\texttt{#1}}

\newtheoremstyle{claimstyle}
  {6pt}{4pt}{\itshape}{}{\bfseries}{.}{0.5em}{}
\theoremstyle{claimstyle}
\newtheorem{claim}{Claim}[section]
\newtheorem{remark}{Remark}[section]

\titleformat{\section}{\large\bfseries}{\thesection.}{0.6em}{}
\titleformat{\subsection}{\normalsize\bfseries}{\thesubsection}{0.6em}{}
\titleformat{\subsubsection}{\normalsize\itshape}{\thesubsubsection}{0.6em}{}

\fancyhf{}
\fancyhead[R]{\small\thepage}
\renewcommand{\headrulewidth}{0.4pt}

\renewcommand{\abstractnamefont}{\normalfont\bfseries}
\renewcommand{\abstracttextfont}{\normalfont\small}
\setlength{\absleftindent}{1.5em}
\setlength{\absrightindent}{1.5em}

\DeclareMathOperator{\resolve}{resolve}

\begin{document}
\author{Pedro Martins Rodrigues Novaes}
\begin{center}
{\LARGE\bfseries Cascading Discreet Log Contracts}\\[0.7em]
{\large Pedro Martins Rodrigues Novaes, GPT 5.5}
\end{center}

\begin{abstract}
Cascading Discreet Log Contracts, or cDLCs, are a construction for composing
DLCs into finite graphs of conditionally activated Bitcoin transactions. For
each parent outcome \(x\), the oracle attestation point \(S_x\) is used as the
adaptor point for signatures on a bridge transaction that spends the parent
outcome output into a prepared child funding output. When the oracle later
reveals the attestation scalar \(s_x\), the same scalar that settles the parent
outcome also completes the bridge signatures and funds the child DLC.
\end{abstract}

\tableofcontents
\newpage

%-----------------------------------------------------------------------
\section{Notation}
%-----------------------------------------------------------------------

Let $G$ be the generator of the secp256k1 group of order $n$. Lowercase
letters denote scalars modulo $n$; uppercase letters denote curve points. For a
scalar $x$, write $X = xG$. Let $H$ denote a tagged hash interpreted as a
scalar modulo $n$.

A BIP340 Schnorr signature for message $m$ under public key $P = pG$ is a pair
$(R,s)$ satisfying
\[
  sG = R + H(R \,\|\, P \,\|\, m)\,P.
\]

An oracle has signing key $v$, public key $V = vG$, and commits to a one-time
nonce point $R_o = r_o G$ for an event. For outcome $x$, define
\begin{align*}
  e_x &= H(R_o \,\|\, V \,\|\, x),\\
  s_x &= r_o + e_x v \pmod{n},\\
  S_x &= s_x G = R_o + e_x V.
\end{align*}

Before the event, anyone can compute $S_x$ for each outcome $x$, but only the
oracle knows $s_x$. After the event, the oracle attests by publishing $s_x$.
This scalar is the transferable fact that makes DLCs possible.

%-----------------------------------------------------------------------
\section{Signature Adaptors}
%-----------------------------------------------------------------------

Let $T = tG$ be a public point whose discrete logarithm $t$ is unknown. To
create an adaptor signature for message $m$ under key $P = pG$, the signer
chooses nonce $r$ and computes
\begin{align*}
  R   &= rG,\quad R^* = R + T,\\
  e   &= H(R^* \,\|\, P \,\|\, m),\\
  \hat{s} &= r + ep \pmod{n}.
\end{align*}
The adaptor $\bigl(R^*,\, T,\, \hat{s}\bigr)$ can be verified without knowing
$t$:
\[
  \hat{s}\,G = R^* - T + eP.
\]
It is not yet a valid Schnorr signature. Once $t$ is revealed, anyone holding
the adaptor computes $s = \hat{s} + t \pmod{n}$ and $(R^*, s)$ is valid because
\[
  sG = (\hat{s}+t)G = (R^*-T+eP)+T = R^* + eP.
\]
Conversely, a completed signature reveals the hidden scalar to any holder of the
adaptor:
\[
  t = s - \hat{s} \pmod{n}.
\]

%-----------------------------------------------------------------------
\section{A DLC as a Conditional Signature Source}
%-----------------------------------------------------------------------

In a regular DLC, Alice and Bob lock funds into a funding output and prepare a
set of Contract Execution Transactions (CETs). Each CET corresponds to an
outcome $x$ and pays according to a payout function $\pi(x)$.

For each outcome $x$, the parties use the oracle attestation point $S_x$ as the
adaptor point for the CET signatures. The CET for $x$ cannot be completed until
the oracle publishes $s_x$.

Thus a DLC does more than choose a payout: it emits a scalar
\[
  \resolve(\mathrm{DLC},\, x) \;\longrightarrow\; s_x.
\]
This scalar is public after resolution, unpredictable before, and bound to a
specific oracle event and outcome. A cDLC uses this scalar as the activation
secret for the next transaction.

%-----------------------------------------------------------------------
\section{cDLC Construction}
%-----------------------------------------------------------------------

A cDLC is a finite directed graph $\Gamma = (N, E)$ where each node $C_i \in N$
is a DLC and each directed edge
\[
  e = (C_i,\, x,\, C_j)
\]
means: if contract $C_i$ resolves to outcome $x$, activate contract $C_j$.

For each edge $e = (C_i, x, C_j)$ the parties construct a \emph{bridge
transaction} $B_e$. This transaction spends a designated output of the parent
CET $\mathrm{CET}_{i,x}$ and creates the funding output $F_j$ of the child DLC
$C_j$. The parent CET contains an \emph{edge output} $O_e$ spendable by the
required cDLC participants, with a timelocked refund path.

The child contract $C_j$ is negotiated before the parent resolves. Its CETs and
refund transaction spend $F_j$, which is possible because the \code{txid} of
$B_e$ is known before its witness signatures are completed.

All signatures on $B_e$ are adaptor signatures using $T_e = S_{i,x}$.
For participant $a$ with key $P_a = p_a G$, let $m_e$ be the sighash for $B_e$.
The participant computes
\begin{align*}
  R_a   &= r_a G,\quad R^*_a = R_a + S_{i,x},\\
  e_a   &= H(R^*_a \,\|\, P_a \,\|\, m_e),\\
  \hat{s}_a &= r_a + e_a p_a \pmod{n}.
\end{align*}
The adaptor verifies as $\hat{s}_a G = R^*_a - S_{i,x} + e_a P_a$.
Once the oracle publishes $s_{i,x}$, anyone holding the adaptors completes:
\[
  s_a = \hat{s}_a + s_{i,x} \pmod{n},
\]
yielding $s_a G = R^*_a + e_a P_a$, a valid Schnorr signature.

Every adaptor on $B_e$ must be exchanged and verified before
$\mathrm{Funding}_i$ is broadcast. The completability step requires at least one
honest party to retain the adaptor set; the oracle scalar alone is insufficient
if all adaptor holders are unavailable.

%-----------------------------------------------------------------------
\section{Native Chaining}
%-----------------------------------------------------------------------

The bridge transaction spends an output of a \emph{future} CET. This is
possible because the \code{txid} of a SegWit or Taproot transaction does not
include witness data: the parent CET can be fully specified before its witness
signatures are completed.

Thus the bridge and the child DLC are constructed before the parent CET is
broadcast:
\[
  \mathrm{Funding}_i \;\to\; \mathrm{CET}_{i,x} \;\to\; B_e \;\to\; \mathrm{Funding}_j.
\]
The off-chain cascade reads:
\begin{align*}
  &\text{oracle signs } x
  \;\Rightarrow\; s_{i,x} \text{ revealed}
  \;\Rightarrow\; \mathrm{CET}_{i,x} \text{ executable}\\
  &\Rightarrow\; B_e \text{ completable}
  \;\Rightarrow\; C_j \text{ funded.}
\end{align*}
Bitcoin validates only ordinary signatures and timelocks. The chain of meaning
is entirely off-chain.

%-----------------------------------------------------------------------
\section{Security Claims}
%-----------------------------------------------------------------------

\begin{claim}[Conditional Activation]
For an edge $e = (C_i, x, C_j)$, the bridge $B_e$ cannot be completed before
$s_{i,x}$ is known, except by forging a Schnorr signature or solving the
discrete logarithm of $S_{i,x}$.
\end{claim}
\begin{proof}
Each required signature on $B_e$ is an adaptor with point $S_{i,x}$.
Completing it without $s_{i,x}$ requires producing a valid Schnorr signature
for $m_e$ under $P_a$ (contradicting unforgeability) or finding $\log_G
S_{i,x}$.
\end{proof}

\begin{claim}[Public Completion After Resolution]
If the oracle publishes $s_{i,x}$, any party holding the bridge adaptors can
complete $B_e$.
\end{claim}
\begin{proof}
For every required signature: $s_a = \hat{s}_a + s_{i,x}$. Since $S_{i,x} =
s_{i,x}G$, we have $s_a G = (\hat{s}_a + s_{i,x})G = R^*_a + e_a P_a$, a valid
Schnorr signature.
\end{proof}

\begin{claim}[Outcome Isolation]
If the oracle signs outcome $y \neq x$ under the same nonce $R_o$, the scalar
$s_{i,y}$ does not complete the bridge for edge $(C_i, x, C_j)$.
\end{claim}
\begin{proof}
Completing with $s_{i,y}$ yields $(\hat{s}_a + s_{i,y})G = R^*_a - S_{i,x} +
e_a P_a + S_{i,y}$, which equals $R^*_a + e_a P_a$ only if $S_{i,y} = S_{i,x}$.
With unique oracle nonces and collision-resistant tagged hashes, this occurs
only with negligible probability.
\end{proof}

\begin{remark}[Oracle Equivocation]
Claim~3 is conditional on the oracle using a single nonce per event.

\emph{Same-nonce equivocation}: the oracle signs outcomes $x$ and $y$ with the
same $R_o$. Only one bridge completes; additionally, the oracle's private key
$v$ is recoverable.

\emph{Different-nonce equivocation}: the oracle signs two outcomes using
independent nonce pairs. Claim~3's isolation argument does not apply; both child
edges could become funded simultaneously. In both cases the conflicting
signatures are publicly verifiable evidence of oracle misconduct.
\end{remark}

%-----------------------------------------------------------------------
\section{Machine-Checked Algebra}
%-----------------------------------------------------------------------

The algebra above is modelled in Ada/SPARK and checked with GNATprove. Proof
artifacts are in the repository's \texttt{spark/} directory.

\paragraph{Targets.}
\begin{itemize}[noitemsep]
  \item \code{spark/cdlc\_integer\_proofs.gpr} --- symbolic polynomial
    identities over \code{SPARK.Big\_Integers}.
  \item \code{spark/cdlc\_residue\_proofs.gpr} --- bridge/adaptor identities
    over $\mathbb{Z}/97\mathbb{Z}$ with explicit modular reduction.
  \item \code{spark/cdlc\_proofs.gpr} --- same identities with Ada \code{type
    mod 97}, using ghost lemmas for modular rotation and cancellation.
  \item \code{spark/lazy\_cdlc\_financial\_completeness\_proofs.gpr} ---
    conservative branch-step accounting for computable, collateral-bounded
    financial transitions under Lazy Compression.
\end{itemize}

The cDLC algebra models prove:
\begin{enumerate}[noitemsep]
  \item The oracle scalar maps to the advertised attestation point.
  \item A bridge adaptor signature verifies before completion.
  \item Adding the oracle scalar completes the bridge signature.
  \item A completed signature reveals the hidden scalar by subtraction.
  \item A different oracle scalar does not complete the same bridge signature.
\end{enumerate}
These correspond exactly to $S_x = s_x G$, $\hat{s}G = R^* - S_x + eP$,
$s = \hat{s} + s_x$, $sG = R^* + eP$, $s - \hat{s} = s_x$.

The Lazy financial-completeness model proves the finite branch-step obligations
used by Section~\ref{sec:financial-completeness}: positive and negative
financial flows are conservative redistributions, prepared matching branches
activate, wrong or unprepared branches fall back, equivalent compressed states
share a continuation template, in-window prepared branches are executable, and
terminal payout accounting conserves funded collateral.

All targets cited above completed with no unproved checks and no \code{pragma
Assume} statements, using CVC5, Z3, and Alt-Ergo. The proofs establish
algebraic consistency of these identities and finite accounting predicates;
secp256k1, BIP340, SHA-256, transaction serialisation, computability theory, and
the full production protocol remain outside the machine-checked boundary.

%-----------------------------------------------------------------------
\section{Refunds and Failure Modes}
%-----------------------------------------------------------------------

\subsection{Assumptions}

The construction requires:
\begin{enumerate}[noitemsep]
  \item The oracle uses a unique nonce per event and does not sign conflicting
    outcomes for the same nonce commitment.
  \item The oracle publishes the attestation before the fallback timeout.
  \item At least one party retains pre-signed transactions and adaptor
    signatures for every live edge in the preparation window.
  \item Fee rates are handled by CPFP, anchor outputs, or pre-agreed fee
    reserves.
  \item Timelocks are ordered so that parent settlement, bridge activation, and
    child refunds do not race.
\end{enumerate}

\subsection{Timeout Ordering}

Let $\tau_{\mathrm{attest}}$ be the expected oracle attestation time,
$\tau_{\mathrm{bridge}}$ the bridge activation deadline, and
$\tau_{\mathrm{child}}$ the child refund timeout. The required ordering is:
\[
  \tau_{\mathrm{attest}} < \tau_{\mathrm{bridge}} < \tau_{\mathrm{child}}.
\]

%-----------------------------------------------------------------------
\section{Lazy Graph Preparation}
%-----------------------------------------------------------------------

\label{sec:lazy}

The cDLC activation equation is \emph{edge-local}: preparing bridge $B_e$ for
edge $(C_i, x, C_j)$ depends only on the parent outcome, the bridge transaction,
and the child funding output for that edge. This locality allows lazy
preparation.

Let $K \geq 2$ be a preparation-window depth. At active node $C_i$, the live
window is $W_i = \{C_i, C_{i+1}, \ldots, C_{i+K-1}\}$. Before $C_i$ resolves,
the parties negotiate the next window by exchanging adaptor signatures for
$C_{i+K}$ --- an asynchronous bilateral exchange identical in structure to any
DLC negotiation, requiring no simultaneous presence. If no extension is agreed,
the fallback executes.

\paragraph{State bounds.}
For a non-recombining tree with branching factor $b$, depth $D$, and per-node
weight $P$:
\[
  \mathrm{EagerState}(D) = P\,\frac{b^{D+1}-1}{b-1},\qquad
  \mathrm{LazyState}(K)  = P\,\frac{b^{K}-1}{b-1}.
\]
For fixed $K$, $\mathrm{LazyState}(K)$ is independent of $D$: maximum live
retained state is bounded regardless of total product depth. If reachable states
recombine, let $|\Sigma_h|$ be the distinct states at distance $h$; then
\[
  \mathrm{LiveState}_i(K) \;\leq\; \sum_{h=0}^{K-1}|\Sigma_h|\,P_h.
\]
Recombination, payoff compression, and lazy preparation are compatible and
jointly reduce retained state.

%-----------------------------------------------------------------------
\section{Financial Completeness With Lazy Compression}
%-----------------------------------------------------------------------
\label{sec:financial-completeness}

A financial contract over a finite oracle alphabet $\Sigma$ can be represented
as a state machine. Let
\[
  z = (q, m, a, b, r)
\]
where $q$ is finite control, $m$ is encoded contract memory, $a$ is Alice's
claim, $b$ is Bob's claim, and $r$ is reserve collateral. The funded collateral
is conserved:
\[
  a + b + r = T.
\]

Let
\[
  \delta : Z \times \Sigma \to Z
\]
be a computable next-state function, and let
\[
  \sigma : Z \to \mathrm{Option}(A_{\mathrm{payout}}, B_{\mathrm{payout}})
\]
be a computable terminal-payout function. The contract is collateral-bounded
when every terminal payout satisfies
\[
  A_{\mathrm{payout}} + B_{\mathrm{payout}} \leq T.
\]

For each compressed state class $\kappa(z)$ and oracle outcome $x$, the
compiler prepares one cDLC branch
\[
  B_{\kappa(z),x}: z \to \delta(z,x).
\]
The bridge for that branch is adapted to the oracle attestation point $S_x$.
When the oracle publishes $s_x$, the matching branch becomes completable and
funds the child state. A non-matching outcome scalar does not complete that
branch by outcome isolation.

Positive and negative financial flows are represented as conservative
redistributions of funded sats:
\[
  (a,b,r) \to (a+d,b-d,r),\qquad
  (a,b,r) \to (a-d,b+d,r),\qquad
  (a,b,r) \to (a,b,r).
\]
Each transition is defined only when the paying side has sufficient collateral.

\begin{claim}[Financial Turing Completeness, Bounded Form]
For every computable, collateral-bounded financial contract over oracle streams,
there exists a cDLC family with Lazy Compression that implements the contract's
terminal payoff on every terminating path for which the required preparation
windows and oracle attestations occur.
\end{claim}
\begin{proof}
Let $F$ be a computable, collateral-bounded financial contract. Since $F$ is
computable, there is a finite program computing its transition function
$\delta$ and terminal payout function $\sigma$. Encode the program counter,
working memory, and funded balances as the state $z$ above.

For every compressed state class $\kappa(z)$ and oracle outcome $x$, prepare a
cDLC branch whose bridge transaction maps $z$ to $\delta(z,x)$ and whose
adaptor point is $S_x$. By the cDLC activation equation, publication of $s_x$
completes exactly the matching prepared branch, while wrong outcomes do not
activate it. By construction, the child funding output carries the balance
distribution of $\delta(z,x)$.
\end{proof}

%-----------------------------------------------------------------------
\section{Applications}
%-----------------------------------------------------------------------

A cDLC can express rolling contracts, re-hedging, periodic synthetic exposure,
and conditional refinancing. These products are more precisely described as
\textbf{renewable fixed-term contracts with bilateral extension rights}: either
party may decline to extend the window, causing the fallback to execute.

%-----------------------------------------------------------------------
\section{Lightning Network Extension}
%-----------------------------------------------------------------------

This section shows that when the parties have already prepared a child state,
the parent outcome scalar can also serve as a channel witness. The child state
must still be negotiated in advance; Lightning only carries the conditional
value transfer.

\subsection{Channel Condition Model}

A Lightning channel is modelled as $Q = (A, B, P)$ where $A$ and $B$ are
balances and $P$ is the set of pending conditional transfers. For a pending
transfer $L = (\mathrm{amount}, \mathrm{lock}, \mathrm{expiry})$, a redemption
predicate $\mathrm{redeem}(\mathrm{lock}, w)$ determines settlement. The channel
conserves capacity: $A + B = A' + B'$ in both cases.

\subsection{HTLC Encoding}

For each oracle outcome $x$, define $w_x = \mathrm{enc}(s_x)$ and $h_x =
H_{\mathrm{pay}}(w_x)$. The HTLC lock is $h_x$; the oracle must precommit to
$h_x$ in its event announcement.

\begin{claim}[HTLC Settlement]
If the oracle publishes $s_x$ and $h_x = H_{\mathrm{pay}}(\mathrm{enc}(s_x))$,
the HTLC is redeemable by witness $\mathrm{enc}(s_x)$.
\end{claim}
\begin{proof}
$\mathrm{redeem}(h_x, w) = \mathrm{true}$ iff $H_{\mathrm{pay}}(w) = h_x$.
Substituting $w = \mathrm{enc}(s_x)$ holds by construction.
\end{proof}

\begin{claim}[HTLC Outcome Isolation]
If $H_{\mathrm{pay}}(\mathrm{enc}(s_y)) \neq H_{\mathrm{pay}}(\mathrm{enc}(s_x))$,
then $\mathrm{enc}(s_y)$ does not redeem the HTLC locked to $h_x$.
\end{claim}
\begin{proof}
Acceptance requires $H_{\mathrm{pay}}(\mathrm{enc}(s_y)) = h_x$, contradicting
the premise under collision and second-preimage resistance.
\end{proof}

\begin{claim}[HTLC--Adaptor Synchronisation]
Publication of $s_x$ both redeems the HTLC locked to $h_x$ and completes the
adaptor with point $S_x$.
\end{claim}
\begin{proof}
HTLC redemption follows from Claim~4.1. Adaptor completion: $(\hat{s}+s_x)G =
(R^*-S_x+eP)+S_x = R^*+eP$.
\end{proof}

\subsection{Routed HTLC cDLCs}

For a route $N_0 \to \cdots \to N_k$ with hash lock $h_x$ at every hop, the
standard Lightning expiry ordering applies, with the additional constraint
$\mathrm{expiry}_k > \tau_{\mathrm{attest}} + \delta$ for oracle observation
margin $\delta$.

\begin{claim}[Routed HTLC Atomicity]
A revealed $\mathrm{enc}(s_x)$ propagates settlement back through the route;
without it, the route times out.
\end{claim}
\begin{proof}
For every hop, $\mathrm{redeem}(h_x, \mathrm{enc}(s_x)) = \mathrm{true}$ by
Claim~4.1. Each node that learns the witness from its outgoing channel has the
exact witness for its incoming channel. Without the witness, no hop can satisfy
$H_{\mathrm{pay}}(w) = h_x$ except by preimage or collision.
\end{proof}

\subsection{Point-Locked Encoding}

A point-locked transfer uses lock $S_x$ and witness $s_x$ with predicate
$\mathrm{redeem}_{\mathrm{PTLC}}(S_x, s) = \mathrm{true}$ iff $sG = S_x$.
Since the ordinary DLC announcement already provides $S_x = s_x G$, no
additional oracle hash commitment is required. PTLC deployment requires channel
protocol changes not present in current major Lightning implementations.

\begin{claim}[Point-Lock Settlement]
Oracle publication of $s_x$ redeems a transfer locked to $S_x$, since
$s_x G = S_x$.
\end{claim}

\begin{claim}[Point-Lock Outcome Isolation]
If $s_y G \neq S_x$, then $s_y$ does not redeem the point lock. In the
prime-order group, $s_y G = s_x G$ implies $s_y = s_x \pmod{n}$.
\end{claim}

\subsection{Routed Point-Locked cDLCs}

Let the receiver's point be $T_k = S_x$. For each hop $i$, the node chooses a
private tweak $d_i$ and defines the upstream point $T_{i-1} = T_i + d_i G$.
If downstream settlement reveals $t_i$ with $t_i G = T_i$, the upstream witness
is $t_{i-1} = t_i + d_i \pmod{n}$.

\begin{claim}[Route Tweak Correctness]
$t_{i-1}G = (t_i+d_i)G = T_i + d_i G = T_{i-1}$.
\end{claim}

\subsection{Machine-Checked Lightning Model}

The Lightning extension is modelled in Ada/SPARK, proof target
\code{spark/lightning\_cdlc\_proofs.gpr}, using \code{type mod 97} for finite
scalars and the ideal injective digest $\mathrm{Hash\_Of}(s) = s$. The target
proves 13 companion claims covering HTLC and PTLC settlement, outcome isolation,
route atomicity, tweak correctness, state activation and non-activation, channel
balance conservation, and HTLC--adaptor synchronisation. The run analysed 51
SPARK subprograms and discharged 118 obligations with no unproved checks and no
\code{pragma Assume} statements.

%-----------------------------------------------------------------------
\section{Limitations}
\label{sec:limitations}
%-----------------------------------------------------------------------

The construction is an activation primitive, not a complete production protocol.

\subsection{Cryptographic Assumptions}

The construction assumes discrete logarithm hardness in secp256k1, Schnorr
unforgeability, collision-resistant hash functions, and unique oracle nonces.
Oracle nonce reuse in a cDLC graph compromises the entire downstream graph.

\subsection{Oracle and Liveness Risk}

If the oracle never publishes the outcome scalar, the child graph is not
activated and funds follow the fallback policy. In HTLC mode, a wrong
precommitted $h_x$ causes liveness failure for the channel edge even when the
DLC algebra is satisfied.

\subsection{State and Negotiation Cost}

Lazy preparation bounds maximum live retained state to $\mathrm{LazyState}(K)$,
independent of total product depth $D$. Cumulative negotiation work scales with
$D \cdot b$: long-running contracts require bilateral window extension at every
period. If one party will not negotiate, the fallback executes.

Lazy Compression does not remove the need to prepare the live window before
use. The financial-completeness theorem assumes that the required window has
been materialized and that the oracle publishes the relevant attestations before
the applicable deadlines.

\subsection{Bitcoin Policy, Fees, and Timelocks}

The proofs guarantee that the right scalar completes the right signature, not
that a transaction confirms. Production systems must handle feerates, mempool
policy, package relay, transaction pinning, anchor outputs, CPFP/RBF, dust
limits, and timeout margins.

\subsection{Privacy Limits}

On-chain cDLCs are discreet at the script level under Taproot. In HTLC
Lightning cDLCs, all hops observe the same $h_x$. PTLC cDLC channels resolve
this via per-hop point tweaks.

\subsection{Collateral Requirements}

Every DLC node requires fully collateralized funding with no netting across
positions. Capital locked in the funding output cannot be deployed elsewhere for
the duration of the period.

\subsection{Scope of the Formal Claim}

Under the stated cryptographic assumptions, if all parties have exchanged and
verified adaptors for every live edge in the current preparation window, at
least one party retains the full adaptor state, and the oracle publishes exactly
one valid outcome scalar before the bridge activation deadline, then that scalar
activates the corresponding prepared child edge and cannot activate
non-corresponding edges except through the failure modes stated above.

%-----------------------------------------------------------------------
\section{Conclusion}
%-----------------------------------------------------------------------

A DLC already contains a hidden scalar that becomes public only when a specific
real-world outcome is attested. A cDLC uses that scalar twice: first to settle
the parent contract, and second to complete adaptor signatures that fund the
next contract.

The core equation is:
\[
  S_x = R_o + H(R_o \,\|\, V \,\|\, x)\,V,\qquad
  s_x G = S_x,\qquad
  s = \hat{s} + s_x \pmod{n}.
\]

This is sufficient to build a graph of conditional Bitcoin transactions in
which one DLC outcome activates another.

%-----------------------------------------------------------------------
\begin{thebibliography}{9}
\bibitem{dryja} T.\ Dryja, ``Discreet Log Contracts,'' MIT Digital Currency
  Initiative, \url{https://adiabat.github.io/dlc.pdf}.
\bibitem{bip340} BIP340, ``Schnorr Signatures for secp256k1,''
  \url{https://github.com/bitcoin/bips/blob/master/bip-0340.mediawiki}.
\bibitem{dlcspecs} ``Discreet Log Contract Interoperability Specifications,''
  \url{https://github.com/discreetlogcontracts/dlcspecs}.
\bibitem{adaptors} Bitcoin Optech, ``Adaptor signatures,''
  \url{https://bitcoinops.org/en/topics/adaptor-signatures/}.
\bibitem{bolt2} Lightning BOLT \#2,
  \url{https://github.com/lightning/bolts/blob/master/02-peer-protocol.md}.
\bibitem{bolt3} Lightning BOLT \#3,
  \url{https://github.com/lightning/bolts/blob/master/03-transactions.md}.
\bibitem{bolt4} Lightning BOLT \#4,
  \url{https://github.com/lightning/bolts/blob/master/04-onion-routing.md}.
\bibitem{bolt11} Lightning BOLT \#11,
  \url{https://github.com/lightning/bolts/blob/master/11-payment-encoding.md}.
\bibitem{ptlc} Bitcoin Optech, ``Point Time Locked Contracts,''
  \url{https://bitcoinops.org/en/topics/ptlc/}.
\end{thebibliography}

\end{document}
